\subsection{多项式环中的零因子}

命题7。——设 \(f \in A[X]\) 是单个未定元的非零多项式，且 \(\alpha\) 为其首项系数。如果 \(\alpha\) 在 \(A\) 中可消去（特别地，若 \(f\) 是首一多项式），则对任意非零元素 \(g\) （在 \(A[X]\) 中），我们有

\[
f g \neq 0 \quad a n d \quad \deg (f g) = \deg f + \deg g.
\]

设 \(g \in A[X]\) 是非零多项式， \(\beta\) 为其首项系数， \(n = \deg f\) 且 \(p = \deg g\) 。于是 \(X^{n+p}\) 在 \(fg\) 中的系数是 \(\alpha\beta\) ，它不为零，由此命题得证。

命题8. -- 若 A 是整环，则 \(\mathbf{A}[(\mathbf{X}_i)_{i\in I}]\) 也是整环。

设 \(u, v\) 为 \(\mathbf{A}[(\mathbf{X}_i)_{i \in I}]\) 的两个非零元素；我们必须证明 \(uv \neq 0\) 。现在 \(u\) 与 \(v\) 属于环 \(\mathbf{A}[(\mathbf{X}_j)_{j \in J}]\) ，其中 \(J\) 是 \(I\) 的有限子集。因此我们可以只限于 \(I\) 有限且等于 \(\{1, 2, \ldots, p\}\) 的情形。另一方面，环 \(\mathbf{A}[\mathbf{X}_1, \ldots, \mathbf{X}_p]\) 同构于以 \(\mathbf{X}_p\) 为未定元、系数在 \(\mathbf{A}[\mathbf{X}_1, \ldots, \mathbf{X}_{p-1}]\) 中的多项式环。于是，对 \(p\) 作归纳，我们被归结为对 \(\mathbf{A}[\mathbf{X}]\) 证明该命题，而这时只需应用命题 7 即可。

推论1. --- 若 A 是整环，且 u 、 v 是 \(\mathbf{A}[(\mathbf{X}_{i})_{i\in I}]\) 的元素，则 \(\deg(uv)=\deg u+\deg v\) 。

我们可以只限于 \(u\) 与 \(v\) 均非零的情形。设 \(m = \deg u\) ， \(n = \deg v\) ；我们有

\[
u = u _ {0} + u _ {1} + \dots + u _ {m}, v = v _ {0} + v _ {1} + \dots + v _ {n}
\]

，其中 \(u_{h}\) （相应地 \(v_{h}\) ）是 u（相应地 v）的 h 次齐次分量。由于 \(u_{m} \neq 0\) 且 \(v_{n} \neq 0\) ，我们有 \(u_{m}v_{n} \neq 0\) （命题 8）。现在 \(uv = u_{m}v_{n} + w\) ，其中 \(\deg w < m + n\) ，由此得出结果。

推论2.——若 A 是整环，则 A{[}(X \(_{i}\) ) \(_{i \in I}\) {]} 的可逆元就是 A 的可逆元。

这直接由推论1得出。

命题9. —— 设 \(u \in A[(X_i)_{i \in I}]\) ；那么 \(u\) 在环 \(A[(X_i)_{i \in I}]\) 中为幂零的充分必要条件是，它的所有系数在环 \(A\) 中均为幂零。

如同命题 8 的证明，我们可以化归到一元多项式的情形。若 u 的所有系数都是幂零的，则 u 是幂零的（I，第 99 页，推论 1）。假设 \(u\) 幂零但非零，设 \(n\) 为其次数；我们对 \(n\) 作归纳。设 \(\alpha\) 为 \(u\) 的首项系数。存在整数 \(m > 0\) 使得 \(u^m = 0\) 。 \(u^m\) 的首项系数是 \(\alpha^m\) ，因此 \(\alpha^m = 0\) 。现在 \(u - \alpha X^n\) 是幂零的（I，同上引文），且归纳假设表明 \(u - \alpha X^n\) 的所有系数都是幂零的。于是 \(u\) 的所有系数都是幂零的。

注记。--- 设 u 和 v 为 \(\mathbf{A}[(\mathbf{X}_{i})_{i\in I}]\) 的元素，并假设 A 是整环， v 是 u 的非零倍元且 v 是齐次的；则 u 也是齐次的。事实上，设 \(u'\in\mathbf{A}[(\mathbf{X}_{i})_{i\in I}]\) 使得 \(v=uu'\) ；我们有 \(u\neq0\) ， \(u'\neq0\) ，且若

\[
\begin{array}{c} u = u _ {h} + u _ {h + 1} + \dots + u _ {k} \\ u ^ {\prime} = u _ {h ^ {\prime}} ^ {\prime} + u _ {h ^ {\prime} + 1} ^ {\prime} + \dots + u _ {k ^ {\prime}} ^ {\prime} \end{array}
\]

是 u 与 \(u'\) 的齐次分量分解，其中 \(u_{h} \neq 0\) , \(u_{k} \neq 0\) , \(u_{h'}' \neq 0\) , \(u_{k'}' \neq 0\) ，则 \(v = u_{h} u_{h'}' + u_{h} u_{h'+1}' + \cdots + u_{k} u_{k'}'\) ，且 \(u_{h} u_{h'}'\) 是次数为 \(h + h'\) 的非零齐次元，而 \(u_{k} u_{k'}'\) 是次数为 \(k + k'\) 的非零齐次元（命题 8）。由于 v 是齐次的，我们有 \(h + h' = k + k'\) ，从而 h = k ， \(h' = k'\) 。

